% !TeX root = ./thuthesis-example.tex

% 论文基本信息配置

% 技术方案基本信息配置（简化版）

% 修复 TimesNewRoman 不支持小型大写字母 (small caps) 的警告
% fontspec 加载 Times New Roman 后内部族名为 TimesNewRoman(0)
% 声明 sc 字形回退到 n，消除 NFSS undefined shape 警告
\AtEndPreamble{%
  \DeclareFontShape{TU}{TimesNewRoman(0)}{m}{sc}{<->ssub*TimesNewRoman(0)/m/n}{}%
  \DeclareFontShape{TU}{TimesNewRoman(0)}{bx}{sc}{<->ssub*TimesNewRoman(0)/bx/n}{}%
}

% 载入所需的宏包

% 定理类环境宏包
\usepackage{amsthm}
% 也可以使用 ntheorem
% \usepackage[amsmath,thmmarks,hyperref]{ntheorem}

% 数学字体配置已简化，使用默认设置

% 可以使用 nomencl 生成符号和缩略语说明
% \usepackage{nomencl}
% \makenomenclature

% 表格加脚注
\usepackage{threeparttable}

% 表格中支持跨行
\usepackage{multirow}

% 固定宽度的表格。
% \usepackage{tabularx}

% 跨页表格
\usepackage{longtable}
\usepackage{tabularx}
\usepackage{booktabs}  % 提供了 \toprule, \midrule, \bottomrule
% 算法
\usepackage{algorithm}
\usepackage{algpseudocode}
\usepackage{listings}
%代码字体
\DeclareFixedFont{\ttb}{T1}{txtt}{bx}{n}{10} % for bold
\DeclareFixedFont{\ttm}{T1}{txtt}{m}{n}{10}  % for normal

%代码颜色
\usepackage{color}
\definecolor{deepblue}{rgb}{0,0,0.5}
\definecolor{deepred}{rgb}{0.6,0,0}
\definecolor{deepgreen}{rgb}{0,0.5,0}

% 量和单位
\usepackage{siunitx}

% 参考文献使用 BibTeX + natbib 宏包
% 顺序编码制
\usepackage[sort]{natbib}
\bibliographystyle{thuthesis-numeric}
\usepackage{titlesec}
\usepackage{tikz}

% 让 \paragraph 参与编号（段级别=4）
\setcounter{secnumdepth}{4}

% 定义 \paragraph 的编号样式为纯阿拉伯数字：1, 2, 3, ...
\makeatletter
\renewcommand\theparagraph{\arabic{paragraph}}
% 如果希望每个 section 内从 1 重新开始编号，保留下一行；
% 如果希望全文连续编号，请删掉下一行。
\@addtoreset{paragraph}{section}
\makeatother

% 把 \paragraph 改成块级标题（标题后换行）
\titleformat{\paragraph}[block]
{\normalfont\normalsize\bfseries}% 样式
{\theparagraph}%                  % 显示的编号（比如 1, 2, 3）
{0.8em}%                          % 编号与标题的间距
{}                                % 标题前置代码

% 调整段前段后间距（可按需修改）
\titlespacing*{\paragraph}
{0pt}{3.25ex plus 1ex minus .2ex}{1ex plus .2ex}
\usepackage{indentfirst}  % 首行缩进

% 著者-出版年制
% \usepackage{natbib}
% \bibliographystyle{thuthesis-author-year}

% 生命科学学院要求使用 Cell 参考文献格式（2023 年以前使用 author-date 格式）
% \usepackage{natbib}
% \bibliographystyle{cell}

% 本科生参考文献的著录格式
% \usepackage[sort]{natbib}
% \bibliographystyle{thuthesis-bachelor}

% 参考文献使用 BibLaTeX 宏包
% \usepackage[style=thuthesis-numeric]{biblatex}
% \usepackage[style=thuthesis-author-year]{biblatex}
% \usepackage[style=gb7714-2015]{biblatex}
% \usepackage[style=apa]{biblatex}
% \usepackage[style=mla-new]{biblatex}
% 声明 BibLaTeX 的数据库
% \addbibresource{ref/refs.bib}
%可固定下划线长度
\makeatletter
\newcommand\dlmu[2][4cm]{\hskip1pt\underline{\hb@xt@ #1{\hss#2\hss}}\hskip3pt}

% 按课程模板固定页眉与摘要标题块
\newcommand\course@report@header{计算机信息工程学院大作业报告}
\newcommand\course@report@title{}
\newcommand\course@report@title@en{}
\newcommand\course@report@headformat[1]{%
  \kaishu\fontsize{10.5bp}{12bp}\selectfont #1%
}
\newcommand\course@report@titleformat[1]{%
  \vspace*{18bp}
  \begin{center}
    {\songti\fontsize{16bp}{20bp}\selectfont #1\par}
    \vspace{18bp}
  \end{center}
}
\newcommand\course@report@abstractname[1]{%
  \begin{center}
    {\songti\bfseries\fontsize{12bp}{20bp}\selectfont #1\par}
    \vspace{20bp}
  \end{center}
}
\fancypagestyle{plain}{%
  \fancyhf{}%
  \renewcommand\headrulewidth{0.75bp}%
  \renewcommand\footrulewidth{0pt}%
  \fancyhead[C]{\course@report@headformat{\course@report@header}}%
  \fancyfoot[C]{\wuhao\thepage}%
}
\fancypagestyle{chapter}{%
  \fancyhf{}%
  \renewcommand\headrulewidth{0.75bp}%
  \renewcommand\footrulewidth{0pt}%
  \fancyhead[C]{\course@report@headformat{\course@report@header}}%
  \fancyfoot[C]{\wuhao\thepage}%
}
\fancypagestyle{course@english@abstract}{%
  \fancyhf{}%
  \renewcommand\headrulewidth{0.75bp}%
  \renewcommand\footrulewidth{0pt}%
  \fancyhead[C]{\course@report@headformat{Abstract}}%
  \fancyfoot[C]{\wuhao\thepage}%
}
\fancypagestyle{course@abstract}{%
  \fancyhf{}%
  \renewcommand\headrulewidth{0.75bp}%
  \renewcommand\footrulewidth{0pt}%
  \fancyhead[C]{\course@report@headformat{\course@report@header}}%
}
\fancypagestyle{course@english@abstract@nopage}{%
  \fancyhf{}%
  \renewcommand\headrulewidth{0.75bp}%
  \renewcommand\footrulewidth{0pt}%
  \fancyhead[C]{\course@report@headformat{Abstract}}%
}
\pagestyle{plain}
\RenewDocumentEnvironment{abstract}{}{%
  \thusetup{language = chinese}%
  \pagestyle{course@abstract}%
  \if@openright\cleardoublepage\else\clearpage\fi
  \thu@pdfbookmark{0}{\thu@abstract@name}%
  \course@report@titleformat{\course@report@title}%
  \course@report@abstractname{\thu@abstract@name}%
  \thispagestyle{course@abstract}%
}{%
  \par
  \null\par
  \noindent
  \textsf{关键词：}%
  \thu@clist@use{\thu@keywords}{；}\par
  \gdef\thu@keywords{}%
  \thispagestyle{course@abstract}%
  \thu@reset@main@language
}
\RenewDocumentEnvironment{abstract*}{}{%
  \thusetup{language = english}%
  \pagestyle{course@english@abstract@nopage}%
  \if@openright\cleardoublepage\else\clearpage\fi
  \thu@pdfbookmark{0}{\thu@abstract@name@en}%
  \course@report@titleformat{\course@report@title@en}%
  \course@report@abstractname{\thu@abstract@name@en}%
  \thispagestyle{course@english@abstract@nopage}%
}{%
  \par
  \null\par
  \noindent
  \textbf{Keywords:} \thu@clist@use{\thu@keywords@en}{; }\par
  \thispagestyle{course@english@abstract@nopage}%
  \thu@reset@main@language
}
\makeatother

\titleformat{\section}[block]{\Large\bfseries}{\thesection}{1em}{}
% 定义所有的图片文件在 figures 子目录下
\graphicspath{{figures/}}

% 数学命令
\makeatletter
\newcommand\dif{%  % 微分符号
  \mathop{}\!%
  \ifthu@math@style@TeX
    d%
  \else
    \mathrm{d}%
  \fi
}
\makeatother

\newcommand\cppstyle{\lstset{
    literate={-}{{-}}1,                % 让 '-' 正常显示
    language=C++,                      % 设置语言为C++
    basicstyle=\ttm,                   % 基本字体样式
    morekeywords={self},               % 根据需要添加更多关键词
    keywordstyle=\ttb\color{deepblue}, % 关键词样式
    emph={MyClass,__init__},           % 自定义高亮关键词
    emphstyle=\ttb\color{deepred},     % 高亮关键词样式
    stringstyle=\color{deepgreen},     % 字符串样式
    frame=tb,                          % 在代码块上下添加边框
    showstringspaces=false,            % 不显示字符串中的空格
    caption={C++},                     % 设置标题为 "C++"
    captionpos=t                       % 标题位置在顶部
  }}

% 定义C++环境
\lstnewenvironment{cpp}[1][]
{
  \cppstyle
  \lstset{#1}
}
{}

% hyperref 宏包在最后调用
\usepackage{hyperref}
